\begin{table*}[!t]
\centering
\caption{Positioning of the proposed work relative to related edge-cloud offloading research.}
\label{tab:related_taxonomy}
\scriptsize
\begin{tabularx}{\textwidth}{@{}lXXX@{}}
\toprule
Category & Representative focus & Common limitation for this study & How this paper differs \\
\midrule
Optimization-based MEC offloading & Binary offloading, radio-compute allocation, energy and latency tradeoffs \cite{ref17,ref18,ref19} & Often relies on tractable model assumptions and may not adapt to burst/fault state changes. & Uses learned offloading under generated temporal bursts, congestion, queue carryover, and fault windows. \\
DRL offloading & DQN, DDQN, actor-critic, and online scheduling policies \cite{ref03,ref07,ref12} & Many evaluations emphasize average latency or energy without explicit rejection-aware analysis. & Uses reward and reporting that include SLA miss, rejection, edge pressure, and heavy-task routing behavior. \\
Federated edge intelligence & FedAvg, FedProx, and federated DRL across distributed clients \cite{ref02,ref04,ref05,ref06,ref30,ref31} & Federated structure can be shown without strong stress-test evidence or scenario-level diagnostics. & Trains across non-IID edge zones and reports scenario-wise, multi-seed, and ablation evidence. \\
Joint offloading and allocation & Coupled destination selection and resource assignment \cite{ref17,ref18,ref19} & Large joint action spaces can obscure whether gains come from routing or allocation. & Separates Stage-1 binary offloading from Stage-2 residual resource allocation for interpretability. \\
Benchmark and simulation studies & Synthetic tasks, simulated MEC nodes, and replayed state tables \cite{ref24,ref28} & Independent samples can understate contention and race-like load effects. & Uses \dataset{} with fixed schemas, time-correlated stress windows, and precomputed contention pressure. \\
\bottomrule
\end{tabularx}
\end{table*}
